% =============================================================================
% METRICS EXTENSION: NLI-BASED EVALUATION
% Template: IEEEtran.cls 2015/08/26 version V1.8b
% =============================================================================
% This file contains the Methodology and Experiments sections for the 
% metrics-extension component of the AskQE paper.
% Include this in your main .tex file after the appropriate \section{} commands
% =============================================================================

% -----------------------------------------------------------------------------
% METHODOLOGY SECTION
% -----------------------------------------------------------------------------

\subsection{NLI-Based Semantic Evaluation}
\label{sec:nli_methodology}

Beyond lexical overlap metrics (F1, Exact Match, BLEU, chrF) and embedding-based similarity (SBERT), we propose evaluating the semantic consistency between source and backtranslated answers through Natural Language Inference (NLI). This approach frames answer comparison as a textual entailment task: given a source answer $a_s$ and a backtranslated answer $a_{bt}$, we classify their semantic relationship.

\subsubsection{Problem Formulation}

For each question-answer pair generated by the AskQE pipeline, we obtain:
\begin{itemize}
    \item $a_s$: the answer derived from the source sentence
    \item $a_{bt}$: the answer derived from the backtranslated MT output
\end{itemize}

We model the semantic relationship between these answers as a three-class classification problem:
\begin{equation}
    f(a_s, a_{bt}) \rightarrow \{E, N, C\}
\end{equation}
where $E$ denotes \textsc{Entailment} (semantically equivalent), $N$ denotes \textsc{Neutral} (partial or uncertain overlap), and $C$ denotes \textsc{Contradiction} (semantic conflict).

\subsubsection{Classification Approaches}

We implement two complementary approaches to perform this classification:

\paragraph{NLI Classifier}
We employ a fine-tuned DeBERTa-v3-base model~\cite{he2021deberta} trained on the MultiNLI corpus~\cite{williams2018multinli}. Given the concatenated input $[a_s; a_{bt}]$, the model outputs logits $\mathbf{z} \in \mathbb{R}^3$ which are converted to a probability distribution via the softmax function:
\begin{equation}
    P(y = c | a_s, a_{bt}) = \frac{\exp(z_c)}{\sum_{j \in \{E, N, C\}} \exp(z_j)}
\end{equation}
where $c \in \{E, N, C\}$ and $z_c$ is the logit corresponding to class $c$. The predicted label is:
\begin{equation}
    \hat{y}_{NLI} = \arg\max_{c} \, P(y = c | a_s, a_{bt})
\end{equation}

\paragraph{LLM-as-Judge}
As an alternative, we prompt a Large Language Model (Qwen-2.5 3B Instruct) to classify the relationship. The model is prompted with:

\begin{quote}
\small
\texttt{Given two answers to the same question:}\\
\texttt{Answer 1: \{$a_s$\}}\\
\texttt{Answer 2: \{$a_{bt}$\}}\\
\texttt{Classify their relationship as ENTAILMENT, NEUTRAL, or CONTRADICTION.}
\end{quote}

The model generates token probabilities for each label, which we extract and normalize via softmax to obtain $P_{LLM}(y = c | a_s, a_{bt})$.

\subsubsection{Aggregation Strategy}

For each sentence $s_i$ with $n_i$ question-answer pairs, we compute the distribution of predicted labels:
\begin{equation}
    \pi_c^{(i)} = \frac{1}{n_i} \sum_{j=1}^{n_i} \mathbb{1}[\hat{y}_j = c]
\end{equation}
where $\pi_c^{(i)}$ represents the proportion of questions classified as class $c$ for sentence $i$. A translation with high $\pi_E$ (entailment proportion) suggests semantic preservation, while high $\pi_C$ (contradiction proportion) indicates potential translation errors.

\subsubsection{Confidence Estimation}

We also track the confidence of each prediction:
\begin{equation}
    \text{conf}(\hat{y}) = P(y = \hat{y} | a_s, a_{bt})
\end{equation}

This allows us to analyze not only \textit{what} label is predicted, but \textit{how certain} the model is about its prediction. We hypothesize that higher confidence correlates with more reliable predictions.


% -----------------------------------------------------------------------------
% EXPERIMENTS & RESULTS SECTION
% -----------------------------------------------------------------------------

\subsection{Inter-Model Agreement Analysis}
\label{sec:nli_results}

We evaluate the NLI-based metrics on the BioMQM dataset across five language pairs: German (de), Spanish (es), French (fr), Russian (ru), and Chinese (zh-CN). For each language, we process all sentence-level QA pairs through both the NLI classifier and LLM-Judge, comparing their predictions.

\subsubsection{Overall Agreement}

Table~\ref{tab:confusion_matrix} presents the confusion matrix between the two classification approaches across all 23,449 question-answer pairs.

\begin{table}[htbp]
\centering
\caption{Confusion Matrix: LLM-Judge vs NLI Classifier}
\label{tab:confusion_matrix}
\begin{tabular}{l|ccc|c}
\hline
& \multicolumn{3}{c|}{\textbf{NLI Classifier}} & \\
\textbf{LLM-Judge} & ENT & NEU & CON & \textbf{Total} \\
\hline
ENTAILMENT & 9,566 & 539 & 62 & 10,167 \\
NEUTRAL & 3,161 & 5,079 & 525 & 8,765 \\
CONTRADICTION & 416 & 2,385 & 1,716 & 4,517 \\
\hline
\textbf{Total} & 13,143 & 8,003 & 2,303 & 23,449 \\
\hline
\end{tabular}
\end{table}

The overall agreement rate is \textbf{69.77\%} (16,361 out of 23,449 predictions). This moderate agreement suggests that while both approaches capture similar semantic relationships, they exhibit systematic differences in their classification behavior.

\subsubsection{Distribution Bias}

We observe distinct distributional biases between the two models:

\begin{itemize}
    \item \textbf{NLI Classifier} tends toward \textsc{Entailment}: 56.0\% of predictions vs. 43.4\% for LLM-Judge
    \item \textbf{LLM-Judge} assigns more \textsc{Contradiction}: 19.3\% vs. 9.8\% for NLI
    \item \textsc{Neutral} predictions are comparable: 37.4\% (LLM) vs. 34.1\% (NLI)
\end{itemize}

This asymmetry can be attributed to differences in training objectives: the NLI classifier was trained on curated entailment datasets where positive relationships predominate, while the LLM-as-Judge approach may be more sensitive to subtle semantic mismatches.

\subsubsection{Language-Specific Agreement}

Table~\ref{tab:language_agreement} shows the agreement rate varying across languages.

\begin{table}[htbp]
\centering
\caption{Inter-Model Agreement by Language}
\label{tab:language_agreement}
\begin{tabular}{lccc}
\hline
\textbf{Language} & \textbf{Total} & \textbf{Matches} & \textbf{Agreement} \\
\hline
German (de) & 5,843 & 4,348 & 74.41\% \\
French (fr) & 3,155 & 2,189 & 69.38\% \\
Spanish (es) & 3,995 & 2,738 & 68.54\% \\
Chinese (zh-CN) & 7,459 & 5,099 & 68.36\% \\
Russian (ru) & 2,997 & 1,987 & 66.30\% \\
\hline
\end{tabular}
\end{table}

German exhibits the highest agreement (74.41\%), possibly due to its structural similarity to English, on which both models were primarily trained. Russian shows the lowest agreement (66.30\%), suggesting that morphologically rich languages with different writing systems pose greater challenges for cross-model consistency.

\subsubsection{Confidence Analysis}

A key finding emerges when analyzing prediction confidence. Table~\ref{tab:confidence} compares the average probability assigned to each label under two conditions: (1) when a model independently assigns that label, and (2) when both models agree on the same label.

\begin{table}[htbp]
\centering
\caption{Average Confidence: Independent vs. Agreement}
\label{tab:confidence}
\begin{tabular}{l|cc|cc}
\hline
& \multicolumn{2}{c|}{\textbf{Independent}} & \multicolumn{2}{c}{\textbf{When Agree}} \\
\textbf{Label} & LLM & NLI & LLM & NLI \\
\hline
ENTAILMENT & 96.66\% & 90.91\% & 96.98\% & 92.59\% \\
NEUTRAL & 89.29\% & 88.07\% & 90.11\% & 90.91\% \\
CONTRADICTION & 87.57\% & 81.15\% & 92.40\% & 83.07\% \\
\hline
\end{tabular}
\end{table}

We observe a consistent increase in confidence when both models agree, most notably for \textsc{Contradiction}:

\begin{itemize}
    \item LLM-Judge: $87.57\% \rightarrow 92.40\%$ ($+4.83$ pp)
    \item NLI Classifier: $81.15\% \rightarrow 83.07\%$ ($+1.92$ pp)
\end{itemize}

This finding has practical implications: \textbf{inter-model agreement can serve as a confidence calibration signal}. When both models predict the same label, particularly \textsc{Contradiction}, the prediction is more likely to be correct.

\subsubsection{Implications for Translation Quality Estimation}

The NLI-based approach offers complementary insights to string-matching metrics:

\begin{enumerate}
    \item \textbf{Semantic granularity}: Unlike binary metrics (e.g., Exact Match), NLI provides a three-way classification that distinguishes between complete semantic preservation (ENT), partial overlap (NEU), and semantic conflict (CON).
    
    \item \textbf{Robustness to paraphrasing}: Answers that are semantically equivalent but lexically different (e.g., ``the virus'' vs. ``COVID-19'') would receive low string-matching scores but high entailment probability.
    
    \item \textbf{Ensemble potential}: The 69.77\% agreement rate, combined with increased confidence on agreement, suggests that an ensemble of NLI and LLM-Judge could provide more robust predictions. A simple voting scheme or confidence-weighted combination could leverage the complementary strengths of both approaches.
\end{enumerate}

\subsubsection{Limitations}

Several limitations should be noted:

\begin{itemize}
    \item Both models were primarily trained on English data; performance on low-resource languages may be suboptimal.
    \item The LLM-as-Judge approach is computationally more expensive than the dedicated NLI classifier.
    \item The 30\% disagreement rate indicates non-trivial uncertainty in semantic classification, warranting further investigation into the nature of these disagreements.
\end{itemize}


% -----------------------------------------------------------------------------
% OPTIONAL: ADDITIONAL TABLES FOR APPENDIX
% -----------------------------------------------------------------------------

% If you want to include per-label match statistics in supplementary material:
%
% \begin{table}[htbp]
% \centering
% \caption{Per-Label Agreement Statistics}
% \label{tab:per_label_agreement}
% \begin{tabular}{lcc}
% \hline
% \textbf{Label} & \textbf{Matches} & \textbf{\% of Total} \\
% \hline
% ENTAILMENT & 9,566 & 40.79\% \\
% NEUTRAL & 5,079 & 21.66\% \\
% CONTRADICTION & 1,716 & 7.32\% \\
% \hline
% \textbf{Total} & 16,361 & 69.77\% \\
% \hline
% \end{tabular}
% \end{table}

